\documentclass[12pt,a4paper]{article}
\usepackage{fontspec}
\usepackage{polyglossia}
\setdefaultlanguage{russian}
\setmainfont{STIX Two Text}
\setmonofont{Menlo}
\newfontfamily\cyrillicfonttt{Menlo}
\usepackage[a4paper,margin=2cm]{geometry}
\usepackage{graphicx,longtable,tabularx,array,booktabs}
\usepackage{enumitem,xurl,listings,hyperref}
\hypersetup{colorlinks=true,linkcolor=blue,urlcolor=blue}
\lstset{basicstyle=\ttfamily\footnotesize,breaklines=true,columns=fullflexible,frame=single,keepspaces=true,showstringspaces=false}
\setlength{\parindent}{1.25cm}
\setlength{\parskip}{0.25em}
\renewcommand{\arraystretch}{1.16}
\setlist{nosep,leftmargin=1.5em}
\sloppy
\begin{document}
\begin{titlepage}
\centering
{\large Лабораторная работа № 1\par}
\vspace{2.5cm}
{\LARGE\bfseries Подготовка среды разработки и репозитория защищённого программного проекта\par}
\vspace{1.8cm}
\begin{tabular}{p{4cm}p{11cm}}
Дисциплина: & Основы разработки защищённого программного обеспечения и компьютерных сетей\\
Вариант: & 19 — Система учёта транспортных средств организации\\
Версия сквозного проекта: & v0.1\\
Студент: & Мостовая Полина Андреевна\\
Группа: & ИУСбд-02-23
\end{tabular}
\vfill
{\large 2026\par}
\end{titlepage}
\tableofcontents
\clearpage
\section{Тема и цель работы}

\textbf{Тема:} Подготовка среды разработки и репозитория защищённого программного проекта.

\textbf{Цель работы:} освоить первоначальную настройку среды безопасной разработки, создать структуру программного проекта, настроить виртуальное окружение и Git-репозиторий, а также реализовать минимальную запускаемую версию приложения без размещения конфиденциальных данных в исходном коде.

\section{Индивидуальный вариант}

Номер индивидуального варианта: \textbf{19}.

Название разрабатываемой программной системы: \textbf{«Система учёта транспортных средств организации»}.

Основные роли пользователей системы:

\begin{itemize}
\item \textbf{водитель} — просматривает закреплённые транспортные средства и передаёт сведения о поездках;
\item \textbf{диспетчер} — ведёт учёт транспортных средств, назначает их водителям и контролирует их состояние;
\item \textbf{администратор} — управляет учётными записями, ролями, настройками системы и событиями безопасности.
\end{itemize}
\section{Описание выбранной программной системы}

Разрабатываемая система предназначена для учёта транспортных средств организации. Система будет вести карточки транспортных средств, хранить сведения об их техническом состоянии и закреплении за водителями, а также поддерживать планирование и контроль поездок. К основным функциям относятся:

\begin{itemize}
\item регистрация, аутентификация пользователей и ролевое разграничение доступа;
\item ведение реестра транспортных средств (государственный номер, модель, статус, технические данные);
\item закрепление транспортных средств за водителями и регистрация поездок;
\item поиск и фильтрация сведений о транспортных средствах и поездках;
\item журналирование событий безопасности;
\item формирование учётных отчётов.
\end{itemize}
Обрабатываемыми данными являются идентификационные данные и учётные записи пользователей, сведения о транспортных средствах, поездках и статусах, записи журнала аудита. К защищаемым активам относятся учётные записи пользователей, секретные ключи, прикладные данные, исходный код, журналы безопасности и база данных.

На данном этапе (v0.1) реализуется только первоначальная безопасная структура: загрузка конфигурации из переменных окружения, проверка секретного параметра и минимальная точка запуска приложения.

\section{Перечень использованного программного обеспечения}

\small\begin{tabularx}{\linewidth}{|p{0.8cm}|p{3.8cm}|p{1.6cm}|X|}
\hline
№ & Программное обеспечение & Версия & Назначение \\ \hline
1 & macOS & 26.6.2 & операционная система \\ \hline
2 & Python & 3.12.14 & интерпретатор языка программирования в виртуальном окружении проекта \\ \hline
3 & Git & 2.50.1 & система контроля версий \\ \hline
4 & Visual Studio Code / терминал & — & среда разработки \\ \hline
5 & python-dotenv & 1.2.2 & загрузка настроек из файла \texttt{.\allowbreak{}env} \\ \hline
6 & pytest & 9.0.2 & средство автоматического тестирования \\ \hline
\end{tabularx}\normalsize

Проверка установки выполнялась командами \texttt{python --version} и \texttt{git --version}. В системе использовался вызов \texttt{python} (Python 3.12.14).

\section{Структура каталогов проекта}

\begin{lstlisting}
lab1_osn_po_seti/
├── app/
│   ├── __init__.py
│   ├── config.py
│   └── main.py
├── docs/
│   └── project_description.md
├── tests/
│   └── test_main.py
├── .env
├── .env.example
├── .gitignore
├── README.md
└── requirements.txt
\end{lstlisting}

Каталоги создавались в следующем порядке: сначала выполнялась инициализация Git-репозитория (\texttt{git init}), затем создавалось виртуальное окружение и после этого формировалась структура каталогов \texttt{app}, \texttt{docs}, \texttt{tests} с пустыми файлами-заготовками.

\section{Назначение созданных файлов}

\small\begin{tabularx}{\linewidth}{|p{3.8cm}|X|}
\hline
Элемент & Назначение \\ \hline
\texttt{app/\allowbreak{}} & каталог с исходным кодом приложения \\ \hline
\texttt{app/\allowbreak{}\_\_init\_\_.\allowbreak{}py} & файл, делающий каталог \texttt{app} пакетом Python \\ \hline
\texttt{app/\allowbreak{}config.\allowbreak{}py} & загрузка и проверка настроек приложения \\ \hline
\texttt{app/\allowbreak{}main.\allowbreak{}py} & точка запуска приложения \\ \hline
\texttt{tests/\allowbreak{}} & каталог с автоматическими тестами \\ \hline
\texttt{tests/\allowbreak{}test\_main.\allowbreak{}py} & тесты первоначальной версии приложения \\ \hline
\texttt{docs/\allowbreak{}} & каталог проектной документации \\ \hline
\texttt{docs/\allowbreak{}project\_description.\allowbreak{}md} & описание проектируемой системы \\ \hline
\texttt{.\allowbreak{}env} & реальные локальные настройки; не передаётся в Git \\ \hline
\texttt{.\allowbreak{}env.\allowbreak{}example} & безопасный шаблон настроек без реальных секретов \\ \hline
\texttt{.\allowbreak{}gitignore} & правила исключения файлов из Git-репозитория \\ \hline
\texttt{requirements.\allowbreak{}txt} & зафиксированный перечень зависимостей проекта \\ \hline
\texttt{README.\allowbreak{}md} & основная документация проекта \\ \hline
\end{tabularx}\normalsize

\section{Листинги конфигурационного модуля и точки запуска}

\subsection{\texttt{app/\allowbreak{}config.\allowbreak{}py}}

Файл реализует загрузку настроек из переменных окружения и их проверку; библиотека \texttt{python-dotenv} загружает пары ключ–значение из \texttt{.\allowbreak{}env} в окружение процесса [3].

\begin{lstlisting}
"""Загрузка и проверка конфигурации приложения."""

import os
from dataclasses import dataclass

from dotenv import load_dotenv


@dataclass(frozen=True)
class Settings:
    """Проверенные настройки приложения."""

    app_name: str
    app_env: str
    secret_key: str


def load_settings() -> Settings:
    """Загрузить настройки из переменных окружения и проверить их."""
    load_dotenv()

    app_name = os.getenv("APP_NAME", "Система учета транспортных средств").strip()
    app_env = os.getenv("APP_ENV", "development").strip().lower()
    secret_key = os.getenv("APP_SECRET_KEY", "")
    allowed_environments = {"development", "testing", "production"}

    if not app_name:
        raise RuntimeError("Название приложения не может быть пустым")
    if app_env not in allowed_environments:
        raise RuntimeError("Недопустимое значение APP_ENV")
    if len(secret_key) < 32:
        raise RuntimeError("APP_SECRET_KEY должен содержать не менее 32 символов")

    return Settings(
        app_name=app_name,
        app_env=app_env,
        secret_key=secret_key,
    )
\end{lstlisting}

\subsection{\texttt{app/\allowbreak{}main.\allowbreak{}py}}

Файл содержит точку запуска приложения.

\begin{lstlisting}
"""Точка запуска системы учета транспортных средств."""

from app.config import load_settings


def build_startup_message(app_name: str, app_env: str) -> str:
    """Сформировать сообщение без раскрытия конфиденциальных данных."""
    return f"{app_name}: приложение запущено в режиме {app_env}"


def main() -> None:
    """Запустить приложение."""
    settings = load_settings()
    print(build_startup_message(settings.app_name, settings.app_env))


if __name__ == "__main__":
    main()
\end{lstlisting}

Основные особенности реализации:

\begin{itemize}
\item секретный ключ загружается из переменной окружения и не имеет опасного значения по умолчанию;
\item минимальная длина секрета (не менее 32 символов) проверяется при запуске;
\item допустимые режимы работы (\texttt{development}, \texttt{testing}, \texttt{production}) заданы явно;
\item объект настроек объявлен как \texttt{frozen}-датакласс, то есть его значения нельзя изменить после создания;
\item программа не выводит секретный ключ на экран и не записывает его в журнал.
\end{itemize}
\section{Содержимое \texttt{.\allowbreak{}gitignore} и \texttt{.\allowbreak{}env.\allowbreak{}example}}

\subsection{\texttt{.\allowbreak{}gitignore}}

\begin{lstlisting}
# Виртуальное окружение
.venv/
venv/

# Конфиденциальная конфигурация
.env
.env.*
!.env.example

# Файлы Python
__pycache__/
*.py[cod]

# Результаты тестирования
.pytest_cache/
.coverage
htmlcov/

# Журналы
*.log
logs/

# Системные файлы
.DS_Store
Thumbs.db
\end{lstlisting}

Данный файл исключает из репозитория виртуальное окружение, файл реальной конфигурации \texttt{.\allowbreak{}env}, служебные файлы Python, результаты тестирования, журналы и системные файлы. Правило \texttt{!.\allowbreak{}env.\allowbreak{}example} позволяет хранить безопасный шаблон конфигурации.

\subsection{\texttt{.\allowbreak{}env.\allowbreak{}example}}

\begin{lstlisting}
APP_NAME=Система учета транспортных средств
APP_ENV=development
APP_SECRET_KEY=replace_with_random_value_at_least_32_characters
\end{lstlisting}

В шаблоне приведены названия переменных и безопасное демонстрационное значение для \texttt{APP\_SECRET\_KEY}. На его основе был создан рабочий файл \texttt{.\allowbreak{}env}, в который записано сгенерированное командой \texttt{python -c "import secrets; print(secrets.\allowbreak{}token\_urlsafe(32))"} случайное значение. Реальное значение ключа в отчёт не включается и в репозиторий не публикуется.

\section{Результаты запуска приложения}

Приложение запускалось командой:

\begin{lstlisting}
python -m app.main
\end{lstlisting}

Полученный результат:

\begin{lstlisting}
Система учета транспортных средств: приложение запущено в режиме development
\end{lstlisting}

Приложение успешно стартовало: конфигурация была загружена из файла \texttt{.\allowbreak{}env}, секретный ключ корректно проверен (длина не менее 32 символов), стартовое сообщение сформировано без вывода секретных данных.

\section{Результаты выполнения тестов}

Тесты запускались командой (pytest обнаруживает тесты по принятой структуре имён файлов и функций [4]):

\begin{lstlisting}
python -m pytest -v
\end{lstlisting}

Результат выполнения:

\begin{lstlisting}
tests/test_main.py::test_startup_message_contains_application_name PASSED [ 33%]
tests/test_main.py::test_startup_message_contains_environment PASSED     [ 66%]
tests/test_main.py::test_startup_message_does_not_contain_secret PASSED  [100%]

============================== 3 passed in 0.00s ===============================
\end{lstlisting}

Все три теста успешно пройдены (\texttt{3 passed}). Тесты проверяют, что стартовое сообщение содержит название приложения и режим работы, а также не содержит секретных значений.

\section{Результаты проверки \texttt{git status}}

После добавления разрешённых файлов в индекс выполнялась команда \texttt{git status}. В выводе отображались только подготовленные к фиксации файлы, файл \texttt{.\allowbreak{}env} и каталог \texttt{.\allowbreak{}venv} в списке отсутствовали. Итоговое состояние репозитория:

\begin{lstlisting}
On branch main
Your branch is up to date with 'origin/main'.

nothing to commit, working tree clean
\end{lstlisting}

Дополнительно действие правил \texttt{.\allowbreak{}gitignore} подтверждено командами:

\begin{itemize}
\item \texttt{git check-ignore -v .\allowbreak{}env} $\rightarrow$ правило \texttt{.\allowbreak{}gitignore:6:.\allowbreak{}env};
\item \texttt{git check-ignore -v .\allowbreak{}venv} $\rightarrow$ правило \texttt{.\allowbreak{}gitignore:2:.\allowbreak{}venv/\allowbreak{}};
\item \texttt{git ls-files .\allowbreak{}env} — команда не вывела результата, что подтверждает, что файл \texttt{.\allowbreak{}env} не отслеживается Git.
\end{itemize}
В репозиторий вошли только разрешённые файлы:

\begin{lstlisting}
.env.example
.gitignore
README.md
app/__init__.py
app/config.py
app/main.py
docs/project_description.md
requirements.txt
tests/test_main.py
\end{lstlisting}

\section{История коммитов и метка версии v0.1}

История репозитория (\texttt{git log --oneline}):

\begin{lstlisting}
779fc15 Создана начальная безопасная структура проекта
\end{lstlisting}

Список меток (\texttt{git tag}):

\begin{lstlisting}
v0.1
\end{lstlisting}

Работа выполнялась по стандартной последовательности: изменение файлов $\rightarrow$ добавление в индекс (\texttt{git add .\allowbreak{}}) $\rightarrow$ создание коммита $\rightarrow$ добавление метки версии (\texttt{git tag -a v0.\allowbreak{}1 -m "Начальная версия проекта"}). Коммиту соответствует начальная версия сквозного проекта v0.1.

\section{Проверки ошибочных настроек (самоконтроль)}

Для проверки корректности работы модуля конфигурации выполнены дополнительные проверки.

\textbf{Проверка 1. Отсутствующий секрет.} После временного удаления значения \texttt{APP\_SECRET\_KEY} из файла \texttt{.\allowbreak{}env} приложение завершилось с контролируемой ошибкой \texttt{APP\_SECRET\_KEY должен содержать не менее 32 символов}. Это подтверждает, что проверка длины секрета работает.

\textbf{Проверка 2. Недопустимый режим работы.} При установке \texttt{APP\_ENV=unknown} конфигурация была отклонена приложением с сообщением \texttt{Недопустимое значение APP\_ENV}.

\textbf{Проверка 3. Защита от публикации секрета.} Команда \texttt{git ls-files .\allowbreak{}env} не вывела результата, то есть файл \texttt{.\allowbreak{}env} не отслеживается Git.

После проверок корректные параметры конфигурации были восстановлены.

\section{Анализ реализованных мер безопасности}

\begin{itemize}
\item \textbf{Секреты не хранятся в исходном коде.} Секретный ключ загружается из переменной окружения через файл \texttt{.\allowbreak{}env}, а не прописывается в коде.
\item \textbf{Файл `.env` исключён из Git.} Правила \texttt{.\allowbreak{}gitignore} и команда \texttt{git ls-files} подтверждают, что реальная конфигурация не попадает в репозиторий.
\item \textbf{Проверка конфигурации при запуске.} Приложение проверяет название, режим работы и минимальную длину секретного ключа (не менее 32 символов).
\item \textbf{Отсутствие опасных значений по умолчанию.} Если секрет не задан, используется пустая строка, и приложение корректно завершается.
\item \textbf{Неизменяемость настроек.} Объект \texttt{Settings} объявлен как \texttt{frozen}-датакласс, что исключает случайное изменение параметров после загрузки.
\item \textbf{Секреты не выводятся пользователю.} Стартовое сообщение содержит только название приложения и режим работы.
\item \textbf{Фиксация зависимостей.} Перечень библиотек зафиксирован в \texttt{requirements.\allowbreak{}txt}, что обеспечивает воспроизводимость окружения.
\item \textbf{Изоляция окружения.} Виртуальное окружение \texttt{.\allowbreak{}venv} не входит в репозиторий и может быть восстановлено из \texttt{requirements.\allowbreak{}txt}.
\end{itemize}
Публикация секретного ключа в открытый репозиторий могла бы привести к его компрометации и использованию злоумышленниками для подделки данных или доступа к системам, поэтому вопросу исключения конфиденциальных данных из репозитория в работе уделялось особое внимание.

\section{Ответы на контрольные вопросы}

\begin{enumerate}
\item \textbf{Что включает в себя среда разработки программного обеспечения?}
\end{enumerate}
Среда разработки включает интерпретатор языка программирования, редактор исходного кода (IDE), систему контроля версий, менеджер зависимостей, средства тестирования и статического анализа, конфигурационные файлы и документацию проекта.

\begin{enumerate}
\item \textbf{Для чего применяется система контроля версий Git?}
\end{enumerate}
Git позволяет сохранять историю изменений, сравнивать версии файлов, фиксировать законченные этапы разработки, возвращаться к предыдущим версиям, организовывать совместную работу, а также определять автора и время каждого изменения [1].

\begin{enumerate}
\item \textbf{Чем рабочая директория отличается от индекса и репозитория?}
\end{enumerate}
Рабочая директория — это текущие файлы проекта, с которыми работает разработчик. Индекс — область, в которую файлы помещаются перед фиксацией (подготовленные к коммиту файлы). Репозиторий — хранилище истории сохранённых изменений (коммитов).

\begin{enumerate}
\item \textbf{Для чего создаётся виртуальное окружение Python?}
\end{enumerate}
Виртуальное окружение изолирует зависимости проекта от других проектов и системных пакетов, что предотвращает конфликты версий библиотек, обеспечивает воспроизводимость окружения и упрощает развёртывание [2].

\begin{enumerate}
\item \textbf{Почему каталог `.venv` не следует добавлять в репозиторий?}
\end{enumerate}
Каталог виртуального окружения содержит большое количество служебных файлов, не относящихся к коду проекта, и может быть восстановлен из списка зависимостей \texttt{requirements.\allowbreak{}txt}.

\begin{enumerate}
\item \textbf{Какие сведения относятся к секретам программного проекта?}
\end{enumerate}
К секретам относятся пароли, токены доступа, ключи API, криптографические ключи, строки подключения к базам данных, учётные данные администраторов и закрытые сертификаты.

\begin{enumerate}
\item \textbf{Почему пароли и ключи нельзя хранить непосредственно в исходном коде?}
\end{enumerate}
Исходный код обычно попадает в Git-репозиторий и может быть опубликован или просмотрен другими людьми. Это приводит к утечке конфиденциальных данных и компрометации системы.

\begin{enumerate}
\item \textbf{Чем `.env` отличается от `.env.example`?}
\end{enumerate}
Файл \texttt{.\allowbreak{}env} содержит реальные значения настроек, включая секреты, и хранится только локально. Файл \texttt{.\allowbreak{}env.\allowbreak{}example} содержит названия переменных и безопасные примеры заполнения и включается в репозиторий.

\begin{enumerate}
\item \textbf{Для чего используется файл `.gitignore`?}
\end{enumerate}
\texttt{.\allowbreak{}gitignore} задаёт перечень файлов и каталогов, исключаемых из репозитория: виртуальные окружения, конфиденциальную конфигурацию, служебные и системные файлы.

\begin{enumerate}
\item \textbf{Почему секретные параметры необходимо проверять при запуске?}
\end{enumerate}
Проверка позволяет своевременно обнаружить отсутствующие или некорректные настройки и не допустить запуск приложения с небезопасной конфигурацией, например с коротким или неустановленным секретным ключом.

\begin{enumerate}
\item \textbf{Для чего фиксируются версии программных зависимостей?}
\end{enumerate}
Фиксация версий в \texttt{requirements.\allowbreak{}txt} обеспечивает воспроизводимость окружения: на любой машине будут установлены одинаковые версии библиотек, что исключает ошибки, связанные с несовместимостью версий.

\begin{enumerate}
\item \textbf{Что представляет собой Git-коммит?}
\end{enumerate}
Коммит — это зафиксированное состояние проекта в определённый момент времени. Он сохраняет снимок изменений, автора, дату и сообщение о внесённых изменениях.

\begin{enumerate}
\item \textbf{Для чего используется метка v0.1?}
\end{enumerate}
Метка (тег) \texttt{v0.\allowbreak{}1} отмечает в истории Git конкретную версию проекта — начальную безопасную структуру. В дальнейшем по метке можно будет вернуться к данной версии.

\begin{enumerate}
\item \textbf{Какие файлы первоначальной версии проекта могут быть опубликованы?}
\end{enumerate}
Опубликованы могут быть файлы исходного кода, документация, шаблон конфигурации \texttt{.\allowbreak{}env.\allowbreak{}example}, список зависимостей и правила \texttt{.\allowbreak{}gitignore}. Не подлежат публикации файл \texttt{.\allowbreak{}env} и виртуальное окружение.

\begin{enumerate}
\item \textbf{Какие последствия может вызвать случайная публикация секретного ключа?}
\end{enumerate}
Публикация ключа может позволить злоумышленникам подделывать данные, получать несанкционированный доступ к системам и сервисам, что приведёт к нарушению конфиденциальности, целостности и доступности информации.

\section{Вывод}

В ходе лабораторной работы была создана начальная версия защищённого программного проекта v0.1 «Система учёта транспортных средств организации» (вариант 19). Выполнены следующие задачи:

\begin{itemize}
\item настроено виртуальное окружение Python и установлены зависимости (\texttt{python-dotenv}, \texttt{pytest});
\item сформирована логически организованная структура проекта (\texttt{app}, \texttt{tests}, \texttt{docs});
\item реализовано минимальное запускаемое приложение с загрузкой настроек из переменных окружения и проверкой обязательного секретного параметра;
\item созданы правила \texttt{.\allowbreak{}gitignore}, исключающие конфиденциальные и служебные файлы из репозитория;
\item оформлены \texttt{README.\allowbreak{}md} и \texttt{docs/\allowbreak{}project\_description.\allowbreak{}md};
\item создан автоматический тест, подтвердивший корректность работы приложения (\texttt{3 passed});
\item выполнен первоначальный коммит и зафиксирована метка версии v0.1.
\end{itemize}
Проверки подтвердили, что файл \texttt{.\allowbreak{}env} с реальным секретным ключом не отслеживается Git, а приложение корректно отклоняет некорректную конфигурацию. Освоены основные команды Git и базовые приёмы безопасной разработки. Подготовленная структура будет использоваться в последующих лабораторных работах сквозного проекта.

\clearpage
\section{Список использованных источников}

\begin{enumerate}
\item Git Project. Git Reference [Электронный ресурс]. URL: <https://git-scm.com/docs> (дата обращения: 15.10.2026).
\item Python Software Foundation. \texttt{venv} — Creation of virtual environments. Python 3.12 Documentation [Электронный ресурс]. URL: <https://docs.python.org/3.12/library/venv.html> (дата обращения: 15.10.2026).
\item Saurabh Kumar, Bertrand Bonnefoy-Claudet. python-dotenv [Электронный ресурс]. PyPI. URL: <https://pypi.org/project/python-dotenv/> (дата обращения: 15.10.2026).
\item pytest Development Team. pytest Documentation [Электронный ресурс]. URL: <https://docs.pytest.org/en/stable/> (дата обращения: 15.10.2026).
\end{enumerate}
\end{document}
